\section*{अध्याय \ref{ch:b2:alkene-redox} --- ऐल्कीनों का ऑक्सीकरण और अपचयन}

\begin{solution}{exo:b2:alkene-redox:1}\emergencystretch=3em
मेथिलसाइक्लोहेक्सेन; 1-मेथिलसाइक्लोहेक्सेन-1-ऑल; trans-2-मेथिलसाइक्लोहेक्सेन-1-ऑल (H और B का
सम-योग, अतः \ce{OH} और नया H एक ही ओर, और मेथिल \ce{OH} के trans);
1-मेथिलसाइक्लोहेक्सीन ऑक्साइड; cis-1-मेथिलसाइक्लोहेक्सेन-1,2-डाइऑल; 6-ऑक्सोहेप्टेनल
\ce{CH3CO(CH2)4CHO}।
\end{solution}

\begin{solution}{exo:b2:alkene-redox:2}
हाइड्रोबोरॉनन--ऑक्सीकरण: 2-मेथिलप्रोपेन-1-ऑल। अम्ल-उत्प्रेरित जलयोजन:
2-मेथिलप्रोपेन-2-ऑल।
\end{solution}

\begin{solution}{exo:b2:alkene-redox:3}
trans-2,3-डाइमेथिलऑक्सिरेन, जिसके दोनों मेथिल वलय की विपरीत ओर हैं। यह किरैल है (कोई
दर्पण तल नहीं) और रेसेमिक मिश्रण के रूप में बनता है।
\end{solution}

\begin{solution}{exo:b2:alkene-redox:4}
अपचायी उपचार-पश्चात् प्रक्रिया: प्रोपेनल के दो अणु। ऑक्सीकारी उपचार-पश्चात् प्रक्रिया: प्रोपेनोइक
अम्ल के दो अणु।
\end{solution}

\begin{solution}{exo:b2:alkene-redox:5}
(a) मेसो-ब्यूटेन-2,3-डाइऑल, ध्रुवण-अघूर्णक क्योंकि यह अकिरैल है। (b)
(2\textit{R},3\textit{R}) और (2\textit{S},3\textit{S}) डाइऑल समान मात्रा में, ध्रुवण-अघूर्णक
क्योंकि रेसेमिक।
\end{solution}

\begin{solution}{exo:b2:alkene-redox:6}
अम्ल में जल तृतीयक कार्बन पर आक्रमण करता है, क्षारक में हाइड्रॉक्साइड \ce{CH2} पर; दोनों
स्थितियों में हर कार्बन पर एक \ce{OH} पहुँचता है: वही 2-मेथिलप्रोपेन-1,2-डाइऑल। मेथेनॉल के साथ दोनों
मार्ग भिन्न ईथर देते हैं (जैसा अध्याय में है)।
\end{solution}

\begin{solution}{exo:b2:alkene-redox:7}
विषाक्त पैलेडियम उत्प्रेरक (लिंडलार) पर \ce{H2} के एक तुल्यांक से हाइड्रोजनन:
त्रि-आबंध पर सम-योग (\textit{Z})-ऐल्कीन पर रुक जाता है।
\end{solution}

\begin{solution}{exo:b2:alkene-redox:8}
प्रोपेनोन (3 C) और ब्यूटेनल (4 C) अपने कार्बोनिल कार्बनों पर जुड़ते हैं: \ce{(CH3)2C=CH-CH2CH2CH3},
2-मेथिलहेक्स-2-ईन।
\end{solution}

\begin{solution}{exo:b2:alkene-redox:9}\emergencystretch=3em
\ce{BH3} (या कोई भारी-भरकम बोरेन), फिर \ce{H2O2} और \ce{NaOH}: हेक्सेन-1-ऑल। अम्ल-उत्प्रेरित जलयोजन
द्वितीयक कार्बधनायन से होकर जाता है और हेक्सेन-2-ऑल (मार्कोवनिकोव) देता है, साथ में
पुनर्व्यवस्थित उपोत्पाद।
\end{solution}

\begin{solution}{exo:b2:alkene-redox:10}
\omterm{def:b2:alkene-redox:epoxide}{परॉक्सी अम्ल}, एक इलेक्ट्रॉनरागी, अधिक प्रतिस्थापित, अधिक इलेक्ट्रॉन-समृद्ध
वलय द्वि-आबंध से तेज़ी से अभिक्रिया करता है। भारी-भरकम बोरेन कम बाधित अंतस्थ \ce{C=CH2} से अभिक्रिया करता है।
\end{solution}

\begin{solution}{exo:b2:alkene-redox:11}
trans: \omterm{def:b2:alkene-redox:epoxide}{परॉक्सी अम्ल}, फिर अम्ल में जल (प्रति-विवृति)। cis: किसी पुनःऑक्सीकारक के साथ उत्प्रेरकी \ce{OsO4},
या ठंडा तनु परमैंगनेट (सम)।
\end{solution}

\begin{solution}{exo:b2:alkene-redox:12}
\ce{C6H10} की असंतृप्ति की मात्रा दो है; \ce{H2} के एक तुल्यांक का अर्थ है एक \ce{C=C}, अतः
एक वलय। छह कार्बनों वाला एक अकेला डाइऐल्डिहाइड अपने द्वि-आबंध पर खुले वलय से आता है:
साइक्लोहेक्सीन।
\end{solution}

\begin{solution}{pb:b2:alkene-redox:1}
\textbf{1.} $10 \times 12.011 + 16 \times 1.008 = \qty{136.24}{g/mol}$।
\textbf{2.} $(2 \times 10 + 2 - 16)/2 = 3$।
\textbf{3.} एक वलय और दो द्वि-आबंध।
\textbf{4.} \ce{C10H16 + 2H2 -> C10H20}, 1-आइसोप्रोपिल-4-मेथिलसाइक्लोहेक्सेन।
\textbf{5.} $1.00/136.24 = \qty{7.34}{mmol}$ लिमोनीन, \ce{H2} के \qty{14.68}{mmol}।
\textbf{6.} $V = nRT/p = 0.01468 \times 8.314 \times 298.15/10^5 = \qty{3.64e-4}{m^3}$,
अर्थात् \qty{364}{mL}।
\textbf{7.} नहीं: वलय के कार्बन 1 और 4 में से हर एक पर वलय के दो समान पथ हैं; cis और
trans समावयवी अकिरैल हैं।
\textbf{8.} मेथेनल \ce{HCHO} (1 C), \ce{=CH2} सिरे से, और एक अकेला \ce{C9} यौगिक,
3-ऐसीटिल-6-ऑक्सोहेप्टेनल \ce{CH3CO-CH2CH2-CH(COCH3)-CH2-CHO}।
\textbf{9.} इसके दोनों कार्बन एक ही वलय के हैं: द्वि-आबंध तोड़ने से वलय खुल जाता है
और कुछ भी अलग नहीं होता।
\textbf{10.} मेथेनल (एक गैस)।
\textbf{11.} ऐल्डिहाइड समूह कार्बोक्सिलिक अम्ल बन जाते हैं (मेथेनल आगे \ce{CO2} तक जाता है);
कीटोन अपरिवर्तित रहते हैं।
\textbf{12.} \ce{C9} शृंखला का कीटोन और ऐल्डिहाइड वलय में जुड़े थे; दूसरा
कीटोन (\ce{COCH3}) और मेथेनल पार्श्व शृंखला के द्वि-आबंध के दो सिरे थे।
\textbf{13.} एक (वलय का C4); \omterm{def:b2:alkene-redox:cleavage}{ओज़ोनी-अपघटन} इसे छूता नहीं और यह \ce{C9} उत्पाद के
त्रिविमजनक केंद्र के रूप में बना रहता है।
\textbf{14.} बहिर्वलयी \ce{C=CH2}, जो त्रिप्रतिस्थापित वलय द्वि-आबंध से कम बाधित है।
\textbf{15.} 2-(4-मेथिलसाइक्लोहेक्स-3-ईन-1-इल)प्रोपेन-1-ऑल: \ce{CH2OH} पूर्व अंतस्थ
कार्बन पर।
\textbf{16.} प्राथमिक।
\textbf{17.} \ce{OH} कम प्रतिस्थापित कार्बन पर जाता है, जो जलयोजन के लिए मार्कोवनिकोव के नियम द्वारा
पूर्वानुमानित अभिविन्यास के विपरीत है।
\textbf{18.} नए केंद्र पर दो विन्यास, जो (पहले से मौजूद केंद्र के साथ) दो
अप्रतिबिंबी त्रिविम समावयवी देते हैं।
\textbf{19.} अम्ल-उत्प्रेरित जलयोजन (या कोई अन्य मार्कोवनिकोव जलयोजन): तृतीयक ऐल्कोहॉल,
$\alpha$-टर्पिनिऑल।
\textbf{20.} वलय द्वि-आबंध: त्रिप्रतिस्थापित, अधिक इलेक्ट्रॉन-समृद्ध।
\textbf{21.} 1,2-एपॉक्सी-1-मेथिल-4-(प्रोप-1-ईन-2-इल)साइक्लोहेक्सेन, ऑक्सीजन C1 और C2 के बीच सेतु बनाता हुआ।
\textbf{22.} C1 पर, जो अधिक प्रतिस्थापित कार्बन है (मेथिल धारण करने वाला)।
\textbf{23.} 1-मेथिल-4-(प्रोप-1-ईन-2-इल)साइक्लोहेक्सेन-1,2-डाइऑल, दोनों \ce{OH} trans (प्रति-विवृति)।
\textbf{24.} वलय द्वि-आबंध का \omterm{def:b2:alkene-redox:dihydroxylation}{सम-डाइहाइड्रॉक्सिलन} (उत्प्रेरकी \ce{OsO4}), जो cis
डाइऑल देता है --- यदि अभिकर्मक को उस द्वि-आबंध से वरणात्मक रूप से अभिक्रिया कराई जा सके।
\textbf{25.} लिमोनीन का \qty{1.00}{g} डाइहाइड्रोजन का $\boldsymbol{\approx \qty{364}{mL}}$
अवशोषित करता है।
\end{solution}
